% =====================================================================
%  AI4Good workshop variant.
%  Defines everything main.tex pulls in: title, venue string, abstract,
%  introduction, discussion. Nothing else in the paper changes.
% =====================================================================

\newcommand{\wkTitle}{Portable Causal Fairness Across \\
Synthetic Data Generator Families}

\newcommand{\wkVenue}{AI for Good}

% ---------------------------------------------------------------- abstract
\newcommand{\wkAbstract}{%
Institutions that release synthetic data in place of sensitive records, such as
a hospital, a lending regulator or a statistical agency, choose the process
that generates the table, and can shape that process so unfair pathways are
absent. DECAF made this concrete for one GAN: factor the generator along a
causal graph, sever the edges from a protected attribute to an outcome, and
three fairness definitions become three sets of edge cuts. But DECAF is a
single GAN, carrying no privacy guarantee and no competitive data-quality
story, and the mechanism has never been shown to work in any other
architecture. We ask whether it belongs to DECAF or to causal factorisation
itself, porting all three definitions to nine generators from three unrelated
families (marginals-based, GAN-based and diffusion-based, each with
differentially private variants) across 2{,}520 matched-pair runs on the Adult
and COMPAS datasets. The definitions transfer everywhere: cutting edges
reduces downstream disparity in every family, dataset and privacy level. The
utility cost is small (under two accuracy points  for the cuts worth making in
five of six family--dataset cells \s{this is too specific for the abstract. should simply say, only costs X utility on average (meaning across all experiments,but needn't say this explicitly here). Also, why are you still reporting "accuracy" when you switched it later in the paper to AUC and macro-F1. Be consistent throughout, and in the graphs too!}), and the privacy cost is nil in the
marginals family, because a severed edge is never measured under the privacy
budget. In this setting privacy and fairness turn out to behave close to
independently rather than adversarially. Causal fairness is now available
inside a modern differentially private synthesizer and inside a diffusion
backbone, not only inside one GAN.%
}

% ------------------------------------------------------------------- intro
\newcommand{\wkIntro}{%
Institutions holding sensitive records increasingly publish a synthetic
substitute rather than the records themselves: a hospital, a lending regulator
or a statistical agency fits a generative model and releases a table drawn from
it. The practice is real and growing: the 2020 US Census was released under a
differential-privacy guarantee \citep{abowd2022topdown}, and public challenges
run by NIST, the FDA and the US--UK privacy-technology programme have made
private tabular release a working technology rather than a proposal
\citep{nist2018challenge,V-CHAMPS,vos2024}. That substitution is usually justified on privacy grounds, but it carries a
second, less-exploited consequence. Whoever generates the data chooses the
generating process, and a process can be edited. A disparity that exists in the
world because of a pathway running from a protected attribute, such as national
origin, to an outcome, such as whether a loan is approved, need not exist in
the released table.

The cleanest existing realisation of that lever is
DECAF~\citep{vanbreugel2021decaf}, which factors a generator along a causal
graph so that severing an edge at sampling time makes the generator
structurally unable to use the pathway it named. But DECAF is a single
non-private WGAN-GP~\citep{gulrajani2017wgangp}, whose released tables trail
modern differentially private marginals-based synthesizers and modern
diffusion models on the fidelity and utility benchmarks that decide what
institutions actually deploy \textcolor{red}{(Sec.~\ref{sec:results})}. Whether the fairness mechanism is confined to
that one architecture, or is a property of causal factorisation itself, is
therefore a practical question and not \textcolor{red}{merely} an academic one: it decides whether an
institution deploying synthetic data has the intervention available at all.

We ask two questions of that mechanism. Is edge surgery a property of DECAF,
or of the causal factorisation any such generator uses? And can it be attached
to a generator whose graph is not a supplied causal DAG but the undirected
marginal-dependency graph a differentially private synthesizer builds from
noisy counts? The stakes are practical: if the mechanism that would make a
released table fair is confined to a generator no institution would actually
deploy, it might as well not exist.%
}

% -------------------------------------------------------------- discussion
\newcommand{\wkDiscussion}{%

\textbf{Edge surgery belongs to causal factorisation, not to DECAF.} DECAF
conjectured that its mechanism would extend beyond its own GAN but tested no
other generator. It extends: across nine generators from three unrelated
families, both datasets and every privacy budget, severing edges reduces
downstream disparity, and FTU, DemPar and CF compile to the same edge cuts on
a diffusion backbone and on a differentially private marginals synthesizer as
they do on the original GAN. An institution can now choose a synthesizer for
its fidelity, privacy or maturity and keep the fairness lever, instead of
being pushed onto a non-private GAN to get it.

\textbf{The cut is affordable.} This matters more than the transfer result,
because a fairness mechanism that ruins the release is not a mechanism anyone
will use. It does not ruin it. CF --- the definition worth applying, since it
removes disparity while sparing the pathways a modeller has declared
legitimate --- costs under two accuracy points in five of six
family\,$\times$\,dataset cells, and leaves the released table's downstream
AUC at \CFAUCKEPT\ \s{I don't know what this is...} of its unconstrained value. DemPar, the most aggressive
cut, is the one to think twice about. A publisher is not choosing between a
fair table and a useful one \s{these last two sentences are bad; what was the main take-away about DemPar mechanism? It's fairer, right? I didn't think terrible quality, was it?}.

\textbf{So which generator should a publisher use?} Our results support a
short answer (Figure~\ref{fig:radar}). \emph{Marginals-based synthesizers are
the default} \s{make clear, in the active voice like above, that this was our contribution (except discounting that PreFair did this for MST... Are out results consistent with PreFair?}: they are the truest to the real table on every fidelity measure
we computed, they are differentially private by construction, and the fairness
cut is free in their privacy accounting, because an edge the definition
forbids is simply never measured under the budget --- a property DECAF's GAN
cannot offer, since it pays for every column it draws. \emph{Causal diffusion
is the choice when the fairness lever has to bite}: its fidelity is close
behind, it retains the most downstream signal of any family \s{more than marginals-based}, and the cut
removes far more of the disparity there than on a marginals graph. \emph{The
causal GAN is the one to avoid}: it trails both on fidelity and utility, and
its differentially private variant fails outright. The lever DECAF
introduced is worth having; DECAF is not the place to get it.

\textbf{How much the cut removes is set by the graph, not the generator.} That
last contrast has a simple cause. A hand-supplied causal DAG routes a
protected attribute to the outcome through many edges, leaving a large
disparity to sever; a private marginals graph is \s{not necessarily, but rather, just "tends to be", right?} sparse by construction and
often carries that dependence faintly or not at all. On Adult, FTU on MST is a
literal no-op in all $45$ matched pairs, because the private structure search
never buys the direct \texttt{sex}\,$\rightarrow$\,\texttt{income} edge in the
first place: the constraint is satisfied before it is imposed. So the question
an auditor should ask of a release is not which generator produced it, but
which graph it was sampled along and which attributes were nominated
admissible. \s{this, to me, is a crazy result that raises huge questions, as a reader: If the DAG is supplied for GANs (and for diffusion, for that matter), then why don't we just try supplying a different DAG to reduce cuts (by reducing inadmissable paths) and thereby keeping it strongly connected during sampling, and so increasing fidelity / utility? And by the way, it's not discussed anywhere how the DAG for each role, for each dataset, was chosen!}

\textbf{Privacy and fairness barely interact here.} A natural worry, given
that differential privacy is known to cost the most accuracy on the smallest
groups \citep{bagdasaryan2019disparate}, is that tightening $\varepsilon$
would also blunt any attempt to close a disparity. We do not see that. The cut
removes at least as much disparity under the tight $\varepsilon{=}1$ budget as
under a permissive one. For an institution already publishing under a privacy
guarantee, fairness here is close to a free addition rather than a competing
demand. \s{say this the other way around: "...if already publishing with a fairness guarantee, adding privacy is close to a free addition...".}

\textbf{Only the in-generator cut ships with the table.} Where a private
structure encodes little to sever, as on COMPAS, downstream mitigators remain
the stronger tool \s{I don't know if this is true! AGAIN, I'll restate: I don't think the \textit{reduction} in parity gap is the only measurement to consider; I think parity gap itself is meaningful, right? Like if CF-on-MST is fairer than downstream mitigator, then that's meaningful. Or am I misunderstanding, and is the reduction measuring the same thing, since it is a reduction over the pure-MST vs. mitigated, and pure-MST vs. CF-MST, for example?}. But a public table is used by consumers its publisher will
never meet, and of the four places to intervene the in-generator cut is the
only one that travels with the release and cannot be dropped by whoever uses it, or however it is used. \s{add a final sentence to this para: "so our method is better it makes the fairness built in to the synthetic data from the very start, guaranteeing the fairness throughout its lifetime" or something like that to close with.}

\textbf{Conclusion.} Synthetic data is increasingly the vehicle through which
institutions share what they know for the public good, as when a hospital, a
regulator or a statistical agency releases a table so that others can build on
it. Our results show that the fairness of such a release need not be left to
whoever eventually trains a model on it, nor traded against the privacy and
fidelity these institutions rightly demand: a single causal edit, portable
across the synthesizers they actually deploy, can remove an unfair pathway
before the data ever leaves the building. Making that lever available wherever
data is generated is a small but concrete step toward synthetic releases that
are fair by construction.%
}
